%% 
%% Copyright 2007-2024 Elsevier Ltd
%% 
%% This file is part of the 'Elsarticle Bundle'.
%% ---------------------------------------------
%% 
%% It may be distributed under the conditions of the LaTeX Project Public
%% License, either version 1.3 of this license or (at your option) any
%% later version.  The latest version of this license is in
%%    http://www.latex-project.org/lppl.txt
%% and version 1.3 or later is part of all distributions of LaTeX
%% version 1999/12/01 or later.
%% 
%% The list of all files belonging to the 'Elsarticle Bundle' is
%% given in the file `manifest.txt'.
%% 
%% Template article for Elsevier's document class `elsarticle'
%% with numbered style bibliographic references
%% SP 2008/03/01
%% $Id: elsarticle-template-num.tex 249 2024-04-06 10:51:24Z rishi $
%%
\documentclass[3p,times,twocolumn]{elsarticle}

%% Use the option review to obtain double line spacing
%% \documentclass[authoryear,preprint,review,12pt]{elsarticle}

%% Use the options 1p,twocolumn; 3p; 3p,twocolumn; 5p; or 5p,twocolumn
%% for a journal layout:
%% \documentclass[final,1p,times]{elsarticle}
%% \documentclass[final,1p,times,twocolumn]{elsarticle}
%% \documentclass[final,3p,times]{elsarticle}
%% \documentclass[final,3p,times,twocolumn]{elsarticle}
%% \documentclass[final,5p,times]{elsarticle}
%% \documentclass[final,5p,times,twocolumn]{elsarticle}

%% For including figures, graphicx.sty has been loaded in
%% elsarticle.cls. If you prefer to use the old commands
%% please give \usepackage{epsfig}

%% The amssymb package provides various useful mathematical symbols
\usepackage{amssymb}
%% The amsmath package provides various useful equation environments.
\usepackage{amsmath}
%% The amsthm package provides extended theorem environments
%% \usepackage{amsthm}

\begin{document}

\begin{frontmatter}

%% Title, authors and addresses

%% use the tnoteref command within \title for footnotes;
%% use the tnotetext command for theassociated footnote;
%% use the fnref command within \author or \affiliation for footnotes;
%% use the fntext command for theassociated footnote;
%% use the corref command within \author for corresponding author footnotes;
%% use the cortext command for theassociated footnote;
%% use the ead command for the email address,
%% and the form \ead[url] for the home page:
%% \title{Title\tnoteref{label1}}
%% \tnotetext[label1]{}
%% \author{Name\corref{cor1}\fnref{label2}}
%% \ead{email address}
%% \ead[url]{home page}
%% \fntext[label2]{}
%% \cortext[cor1]{}
%% \affiliation{organization={},
%%             addressline={},
%%             city={},
%%             postcode={},
%%             state={},
%%             country={}}
%% \fntext[label3]{}

\title{IJHCS Template}

%% use optional labels to link authors explicitly to addresses:
%% \author[label1,label2]{}
%% \affiliation[label1]{organization={},
%%             addressline={},
%%             city={},
%%             postcode={},
%%             state={},
%%             country={}}
%%
%% \affiliation[label2]{organization={},
%%             addressline={},
%%             city={},
%%             postcode={},
%%             state={},
%%             country={}}

\author[1]{Author1} %% Author name

%% Author affiliation
\affiliation[1]{organization={},%Department and Organization
            addressline={}, 
            city={},
            postcode={}, 
            state={},
            country={}}

% -----------------------------------------------------------------------------
% THE ABSTRACT
\begin{abstract}
\input{structure/abstract.tex}
\end{abstract}
%
% -----------------------------------------------------------------------------

%% Keywords
\begin{keyword}
% #############################################################################
% Keywords
% !TEX root = main.tex
% #############################################################################

Human-Computer Interaction
\end{keyword}

\end{frontmatter}

%% Add \usepackage{lineno} before \begin{document} and uncomment 
%% following line to enable line numbers
%% \linenumbers

%% main text
%%
% -----------------------------------------------------------------------------
% This a suggestion for the Content of the Document
% Section 1
\input{sections/sec001.tex}
% -----------------------------------------------------------------------------
%%
% -----------------------------------------------------------------------------
% This a suggestion for the Content of the Document
% Section 2
% #############################################################################
% This is Section 2
% !TEX root = main.tex
% #############################################################################
% Change the Name of the Section i the following line
\section{Background}
\label{sec:sec002}
% -----------------------------------------------------------------------------
%%
% -----------------------------------------------------------------------------
% This a suggestion for the Content of the Document
% Section 3
% #############################################################################
% This is Section 3
% !TEX root = main.tex
% #############################################################################
% Change the Name of the Section i the following line
\section{System}
\label{sec:sec003}
% -----------------------------------------------------------------------------
%%
% -----------------------------------------------------------------------------
% This a suggestion for the Content of the Document
% Section 4
% #############################################################################
% This is Section 4
% !TEX root = main.tex
% #############################################################################
% Change the Name of the Section i the following line
\section{Research Questions \& Hypotheses}
\label{sec:sec004}
% -----------------------------------------------------------------------------
%%
% -----------------------------------------------------------------------------
% This a suggestion for the Content of the Document
% Section 5
% #############################################################################
% This is Section 5
% !TEX root = main.tex
% #############################################################################
% Change the Name of the Section i the following line
\section{Methods}
\label{sec:sec005}
% -----------------------------------------------------------------------------
%%
% -----------------------------------------------------------------------------
% This a suggestion for the Content of the Document
% Section 6
% #############################################################################
% This is Section 6
% !TEX root = main.tex
% #############################################################################
% Change the Name of the Section i the following line
\section{Results}
\label{sec:sec006}
% -----------------------------------------------------------------------------
%%
% -----------------------------------------------------------------------------
% This a suggestion for the Content of the Document
% Section 7
% #############################################################################
% This is Section 7
% !TEX root = main.tex
% #############################################################################
% Change the Name of the Section i the following line
\section{Discussion}
\label{sec:sec007}
% -----------------------------------------------------------------------------
%%
% -----------------------------------------------------------------------------
% This a suggestion for the Content of the Document
% Section 8
% #############################################################################
% This is Section 8
% !TEX root = main.tex
% #############################################################################
% Change the Name of the Section i the following line
\section{Conclusion}
\label{sec:sec008}
% -----------------------------------------------------------------------------



%% The Appendices part is started with the command \appendix;
%% appendix sections are then done as normal sections
%%
% -----------------------------------------------------------------------------
% This a suggestion for the Content of the Document
% Section 8
% #############################################################################
% This is Appendix Technical Details
% !TEX root = main.tex
% #############################################################################
\appendix
\section{Example Appendix Section}
\label{app:app001}
% -----------------------------------------------------------------------------

Appendix text.

%% If you have bib database file and want bibtex to generate the
%% bibitems, please use
%%
%%  \bibliographystyle{elsarticle-num} 
%%  \bibliography{<your bibdatabase>}

%% else use the following coding to input the bibitems directly in the
%% TeX file.

%% Refer following link for more details about bibliography and citations.
%% https://en.wikibooks.org/wiki/LaTeX/Bibliography_Management

%% For numbered reference style
%% \bibitem{label}
%% Text of bibliographic item

\end{document}

\endinput
%%
%% End of file `elsarticle-template-num.tex'.
